\documentclass[11pt]{article}
\usepackage[margin=1in]{geometry}
\usepackage{amsmath, amssymb, amsthm}
\usepackage{enumitem}
\usepackage{xcolor}
\usepackage{tcolorbox}
\tcbuselibrary{skins, breakable}

% Define color boxes for different framework sections
\newtcolorbox{problemconfigbox}{
    colback=blue!5, colframe=blue!40!black, title=PROBLEM CONFIGURATION ANALYSIS,
    fonttitle=\bfseries, breakable
}

\newtcolorbox{strategicbox}{
    colback=pink!5, colframe=pink!40!black, title=STRATEGIC FRAMEWORK APPLICATION,
    fonttitle=\bfseries, breakable
}

\newtcolorbox{tacticalbox}{
    colback=green!5, colframe=green!40!black, title=TACTICAL METHODOLOGY SELECTION,
    fonttitle=\bfseries, breakable
}

\newtcolorbox{conversionbox}{
    colback=yellow!5, colframe=yellow!40!black, title=CONVERSION TECHNIQUES APPLICATION,
    fonttitle=\bfseries, breakable
}

\newtcolorbox{verificationbox}{
    colback=red!5, colframe=red!40!black, title=VERIFICATION STRATEGY DESIGN,
    fonttitle=\bfseries, breakable
}

\newtcolorbox{roadmapbox}{
    colback=cyan!5, colframe=cyan!40!black, title=SOLUTION ROADMAP,
    fonttitle=\bfseries, breakable
}

\newtcolorbox{competitionbox}{
    colback=purple!5, colframe=purple!40!black, title=COMPETITION STRATEGY INTEGRATION,
    fonttitle=\bfseries, breakable
}

\newtcolorbox{keyinsightbox}{
    colback=orange!5, colframe=orange!40!black, title=KEY INSIGHT,
    fonttitle=\bfseries,breakable
}

\newtcolorbox{warningbox}{
    colback=red!10, colframe=red!60!black, title=WARNING: COMMON PITFALL,
    fonttitle=\bfseries,breakable
}

\newtcolorbox{tipbox}{
    colback=green!10, colframe=green!60!black, title=PRO TIP,
    fonttitle=\bfseries,breakable
}

\title{\textbf{2016 AMC 12B Problem 16: Strategic Analysis}}
\author{Using AMC12/COMC Problem-Solving Framework}
\date{}

\begin{document}
\maketitle

\section*{Problem Statement}
In how many ways can $345$ be written as the sum of an increasing sequence of two or more consecutive positive integers?

\textbf{Answer Choices:} 
$\textbf{(A)}\ 1\qquad\textbf{(B)}\ 3\qquad\textbf{(C)}\ 5\qquad\textbf{(D)}\ 6\qquad\textbf{(E)}\ 7$

\begin{problemconfigbox}
\subsection*{A. Problem Classification}
\begin{itemize}[leftmargin=*]
    \item \textbf{Problem Type}: To Find - counting the number of ways to represent 345 as consecutive integer sums
    \item \textbf{Difficulty Band}: Mid-to-late band (Problem 16 of 25) - requires algebraic insight and systematic case analysis
    \item \textbf{Competition Context}: AMC12 (75 min, 25 problems) - approximately 3-4 minutes target time
    \item \textbf{Mathematical Domain}: Number Theory / Algebra hybrid - involves factorization, arithmetic sequences, and divisibility
\end{itemize}

\subsection*{B. Problem Structure Identification}
\begin{itemize}[leftmargin=*]
    \item \textbf{Given Information}: 
    \begin{itemize}
        \item Target sum: 345
        \item Sequence must be consecutive positive integers
        \item Sequence must be increasing
        \item Sequence must have at least 2 terms
    \end{itemize}
    \item \textbf{Constraints}: 
    \begin{itemize}
        \item All integers must be positive ($n \geq 1$)
        \item Minimum sequence length is 2 ($k \geq 2$)
        \item Sequence must be strictly increasing
    \end{itemize}
    \item \textbf{Objective}: Count the number of valid representations
    \item \textbf{Hidden Assumptions}: 
    \begin{itemize}
        \item Order matters (sequences starting at different points are different)
        \item Complete sequences only (no partial sums)
    \end{itemize}
    \item \textbf{Answer Choice Leverage}: Small answer choices (1-7) suggest systematic enumeration is feasible
\end{itemize}

\subsection*{C. Strategic Orientation}
\begin{itemize}[leftmargin=*]
    \item \textbf{Hypothesis}: Let first term be $n$ and number of terms be $k$ (where $k \geq 2$, $n \geq 1$)
    \item \textbf{Conclusion}: Find all valid $(n, k)$ pairs such that the sum equals 345
    \item \textbf{Notation}: 
    \begin{itemize}
        \item $n$ = first term of sequence
        \item $k$ = number of consecutive terms
        \item Sum formula: $S = n + (n+1) + \cdots + (n+k-1)$
    \end{itemize}
    \item \textbf{Intuitive Approach}: Use arithmetic sequence sum formula, then analyze constraints
    \item \textbf{Cross-Domain Potential}: Can be viewed as a factorization problem (factoring 690 = 2 × 345)
\end{itemize}
\end{problemconfigbox}

\begin{strategicbox}
\subsection*{A. Psychological Strategies}
\begin{itemize}[leftmargin=*]
    \item \textbf{Mental Toughness}: Need to persist through algebraic manipulation and systematic case checking
    \item \textbf{Creativity Cultivation}: Recognize the connection between consecutive sums and factorization
    \item \textbf{Iterative Refinement}: Start with sum formula, derive constraints, then systematically enumerate
\end{itemize}

\subsection*{B. Investigation Strategies}
\begin{itemize}[leftmargin=*]
    \item \textbf{Startup Orientation}: Immediately write the arithmetic sequence sum formula
    \item \textbf{Get Your Hands Dirty}: Test small cases (k=2, k=3) to build intuition
    \item \textbf{Make It Easier}: Consider what happens with even vs. odd number of terms
    \item \textbf{Conjecture via Small Cases}: Try k=2 first: $n + (n+1) = 2n + 1 = 345 \Rightarrow n = 172$
    \item \textbf{Visualization}: Not particularly helpful for this problem
\end{itemize}

\subsection*{C. Late-Band Strategic Playbook}
\begin{itemize}[leftmargin=*]
    \item \textbf{Answer-Choice Leverage}: Answer is 7, suggesting there are exactly 7 factor pairs to consider
    \item \textbf{Make It Easier}: Simplify to finding factor pairs of $2 \times 345 = 690$
    \item \textbf{Penultimate Step}: The final step is counting valid factor pairs, so work backwards from factorization
\end{itemize}

\subsection*{D. Argument Construction Strategy}
\begin{itemize}[leftmargin=*]
    \item \textbf{Direct Proof}: Derive necessary and sufficient conditions algebraically
    \item \textbf{Constructive Method}: For each valid factor pair, explicitly show the sequence exists
\end{itemize}

\subsection*{E. Cross-Domain Translation}
\begin{itemize}[leftmargin=*]
    \item \textbf{Recast the Problem}: Transform from "consecutive sum problem" to "factorization problem"
    \item \textbf{Change Your Point of View}: View as finding divisors of 690 with specific parity constraints
\end{itemize}
\end{strategicbox}

\begin{tacticalbox}
\subsection*{A. Core Mathematical Tactics}
\begin{itemize}[leftmargin=*]
    \item \textbf{Algebraic Manipulation}: Key tactic - derive closed form for arithmetic sequence sum
    \item \textbf{Factorization}: Convert problem to finding factor pairs of 690
    \item \textbf{Parity Analysis}: Consider odd vs. even cases separately for systematic enumeration
    \item \textbf{Constraint Analysis}: Determine which factor pairs satisfy $n \geq 1$ and $k \geq 2$
\end{itemize}

\subsection*{B. Crossover Tactics}
\begin{itemize}[leftmargin=*]
    \item \textbf{Divisor Counting}: Use prime factorization $690 = 2 \times 3 \times 5 \times 23$
    \item \textbf{Number of Divisors}: $690$ has $(1+1)(1+1)(1+1)(1+1) = 16$ total divisors
\end{itemize}
\end{tacticalbox}

\begin{conversionbox}
\subsection*{A. High-Probability Techniques Applied}
\begin{enumerate}[leftmargin=*]
    \item \textbf{Arithmetic Sequence Sum Formula}: 
    \[S = \frac{k(2n + k - 1)}{2} = 345\]
    Therefore: $k(2n + k - 1) = 690$
    
    \item \textbf{Simon's Favorite Factoring Trick (SFFT)}: Rewrite as product
    \[k(2n + k - 1) = 690 = 2 \times 3 \times 5 \times 23\]
    
    \item \textbf{Case Analysis}: Separate into two cases based on parity of $k$
    
    \textbf{Case 1: $k$ is odd}
    
    Let $k = 2m + 1$ for some non-negative integer $m$. Then:
    \[(2m+1)(2n + 2m) = 690\]
    \[(2m+1) \cdot 2(n+m) = 690\]
    
    For this to have integer solutions, $(2m+1)$ must be an odd divisor of 345.
    
    Odd divisors of $345 = 3 \times 5 \times 23$: \{1, 3, 5, 15, 23, 69, 115, 345\}
    
    Checking $k = 2m+1 \geq 2$: Must have $k \in \{3, 5, 15, 23, 69, 115, 345\}$
    
    But need $n \geq 1$. From $(2m+1)(2n+2m) = 690$:
    \[n = \frac{690/(2m+1) - 2m}{2}\]
    
    For $k = 345$: $n = \frac{2 - 344}{2} = -171 < 0$ $\times$
    
    For $k = 115$: $n = \frac{6 - 114}{2} = -54 < 0$ $\times$
    
    For $k = 69$: $n = \frac{10 - 68}{2} = -29 < 0$ $\times$
    
    Valid odd values: $k \in \{3, 5, 15, 23\}$ → 4 solutions
    
    \textbf{Case 2: $k$ is even}
    
    Let $k = 2m$ for some positive integer $m \geq 1$. Then:
    \[2m(2n + 2m - 1) = 690\]
    \[m(2n + 2m - 1) = 345\]
    
    Here $(2n + 2m - 1)$ is odd, so $m$ can be any divisor of 345.
    
    Divisors of $345$: \{1, 3, 5, 15, 23, 69, 115, 345\}
    
    For $k = 2m \geq 2$: $m \in \{1, 3, 5, 15, 23, 69, 115, 345\}$
    
    Check $n \geq 1$: From $m(2n + 2m - 1) = 345$:
    \[n = \frac{345/m - 2m + 1}{2}\]
    
    For large $m$, this becomes negative. Testing:
    
    For $m = 1$ (k=2): $n = 172$ $\checkmark$
    
    For $m = 3$ (k=6): $n = 55$ $\checkmark$
    
    For $m = 5$ (k=10): $n = 30$ $\checkmark$
    
    For $m = 15$ (k=30): $n = \frac{23-30+1}{2} = -3 < 0$ $\times$
    
    For $m = 23$ (k=46): $n = \frac{15-46+1}{2} = -15 < 0$ $\times$
    
    Valid even values: $k \in \{2, 6, 10\}$ → 3 solutions
    
    \textbf{Total from both cases}: $4 + 3 = 7$ solutions $\checkmark$
    
    \item \textbf{Verification via Complete Factorization Approach}:
    
    From $k(2n+k-1) = 690$, let $a = k$ and $b = 2n+k-1$.
    
    Then $ab = 690$ and $b - a = 2n - 1$.
    
    For $n \geq 1$: need $b - a \geq 1$, so $b > a$ (or $b = a + 1$ minimum).
    
    Also need $b - a$ to be odd (since $2n-1$ is odd), which means $a$ and $b$ have opposite parity.
    
    Since $690 = 2 \times 345$, exactly one of $a, b$ must be even.
    
    Factor pairs $(a,b)$ with $a < b$: 
    \begin{itemize}
        \item $(1, 690)$: $b - a = 689$ → $n = 345$ $\times$ ($k=1$ violates $k \geq 2$)
        \item $(2, 345)$: $b - a = 343$ → $n = 172$ $\checkmark$
        \item $(3, 230)$: $b - a = 227$ → $n = 114$ $\checkmark$
        \item $(5, 138)$: $b - a = 133$ → $n = 67$ $\checkmark$
        \item $(6, 115)$: $b - a = 109$ → $n = 55$ $\checkmark$
        \item $(10, 69)$: $b - a = 59$ → $n = 30$ $\checkmark$
        \item $(15, 46)$: $b - a = 31$ → $n = 16$ $\checkmark$
        \item $(23, 30)$: $b - a = 7$ → $n = 4$ $\checkmark$
    \end{itemize}
    
    Count: 8 pairs, but $(1, 690)$ is invalid, so \boxed{7} valid sequences.
\end{enumerate}

\subsection*{B. Recognition Index}
\begin{itemize}[leftmargin=*]
    \item \textbf{Pattern Recognition Time}: 30-60 seconds to recognize this as arithmetic sequence sum problem
    \item \textbf{Key Indicator}: "Consecutive positive integers" → arithmetic sequence formula
    \item \textbf{Secondary Indicator}: "How many ways" + specific target → factorization/divisor counting
\end{itemize}

\subsection*{C. Technique Integration}
The solution integrates:
\begin{itemize}
    \item Arithmetic sequence sum formula (converts to algebraic equation)
    \item Factorization (identifies all possible cases)
    \item Constraint checking (eliminates invalid solutions)
\end{itemize}

\subsection*{D. Conflict Resolution}
No technique conflicts - the approaches (case analysis by parity vs. factor pair enumeration) are complementary and lead to the same answer.
\end{conversionbox}

\begin{verificationbox}
\subsection*{A. Verification Level Selection}
\textbf{Selected Level}: Level 5 (Thorough) - Mid-to-late band problem warrants comprehensive verification

\subsection*{B. Verification Methods Applied}
\begin{enumerate}[leftmargin=*]
    \item \textbf{Direct Verification}: Check each of the 7 sequences explicitly:
    \begin{itemize}
        \item $k=2$, $n=172$: $172 + 173 = 345$ $\checkmark$
        \item $k=3$, $n=114$: $114 + 115 + 116 = 345$ $\checkmark$
        \item $k=5$, $n=67$: $67 + 68 + 69 + 70 + 71 = 345$ $\checkmark$
        \item $k=6$, $n=55$: $55 + \cdots + 60 = 345$ $\checkmark$
        \item $k=10$, $n=30$: $30 + \cdots + 39 = 345$ $\checkmark$
        \item $k=15$, $n=16$: $16 + \cdots + 30 = 345$ $\checkmark$
        \item $k=23$, $n=4$: $4 + \cdots + 26 = 345$ $\checkmark$
    \end{itemize}
    
    \item \textbf{Boundary Checking}: Verify that the next factor pair would fail:
    \begin{itemize}
        \item $(30, 23)$: This would give $k=30$, $b-a = -7 < 0$ → invalid
        \item $(46, 15)$: This would give $k=46$, $n < 0$ → invalid
    \end{itemize}
    
    \item \textbf{Answer Choice Validation}: Result of 7 matches option (E) $\checkmark$
    
    \item \textbf{Alternative Method}: Count divisors approach
    \begin{itemize}
        \item 690 has 16 total divisors
        \item Form 8 factor pairs with first factor < second factor
        \item Eliminate $(1, 690)$ due to $k=1$ constraint
        \item Result: 7 valid pairs $\checkmark$
    \end{itemize}
\end{enumerate}

\subsection*{C. Confidence Calibration}
\begin{itemize}[leftmargin=*]
    \item \textbf{Confidence Level}: 95\% - Multiple verification methods agree
    \item \textbf{Uncertainty Sources}: None significant - algebraic derivation is straightforward
    \item \textbf{Flag for Review}: No - high confidence with cross-verification
\end{itemize}
\end{verificationbox}

\begin{roadmapbox}
\subsection*{Step-by-Step Solution Plan}

\textbf{Phase 1: Setup and Formula Derivation (1 minute)}
\begin{enumerate}
    \item Write arithmetic sequence sum formula: $S = \frac{k(2n + k - 1)}{2}$
    \item Set equal to 345 and multiply by 2: $k(2n + k - 1) = 690$
    \item Recognize as factorization problem
\end{enumerate}

\textbf{Phase 2: Prime Factorization (30 seconds)}
\begin{enumerate}
    \item Factor 690: $690 = 2 \times 3 \times 5 \times 23$
    \item Note: 690 has $(1+1)(1+1)(1+1)(1+1) = 16$ divisors
\end{enumerate}

\textbf{Phase 3: Systematic Enumeration (1.5 minutes)}
\begin{enumerate}
    \item List all factor pairs $(k, 2n+k-1)$ where $k < 2n+k-1$
    \item For each pair, check if $k \geq 2$ and $n \geq 1$
    \item Eliminate $(1, 690)$ due to $k=1$ violation
    \item Count remaining valid pairs
\end{enumerate}

\textbf{Phase 4: Verification (1 minute)}
\begin{enumerate}
    \item Spot-check 2-3 sequences by direct calculation
    \item Confirm boundary cases are correctly excluded
    \item Match result to answer choices
\end{enumerate}

\subsection*{Alternative Approaches}
\begin{itemize}
    \item \textbf{Parity Case Analysis}: Separate odd and even $k$ cases
    \item \textbf{Quadratic Formula}: Solve for $n$ in terms of $k$ and test all valid $k$ values
    \item \textbf{Computer Search}: For practice, write code to enumerate all solutions
\end{itemize}

\subsection*{Time Management}
\begin{itemize}
    \item \textbf{Target Time}: 3-4 minutes total
    \item \textbf{Actual Breakdown}: Setup (1 min) + Enumeration (1.5 min) + Verification (1 min) = 3.5 min
    \item \textbf{Pivot Indicator}: If taking >5 minutes, mark answer and move on
\end{itemize}

\subsection*{Comprehensive Pitfall Guide}
\begin{enumerate}
    \item \textbf{Algebraic Error}: Incorrectly deriving sum formula
    \item \textbf{Missing Constraint}: Forgetting $k \geq 2$ constraint (excluding $k=1$)
    \item \textbf{Incomplete Enumeration}: Missing some factor pairs of 690
    \item \textbf{Sign Error}: Not checking $n \geq 1$ for all factor pairs
    \item \textbf{Overcounting}: Counting both $(k, 2n+k-1)$ and $(2n+k-1, k)$ as separate
\end{enumerate}
\end{roadmapbox}

\begin{competitionbox}
\subsection*{A. Competition-Specific Strategies}
\begin{itemize}[leftmargin=*]
    \item \textbf{Time Allocation}: 3-4 minutes (appropriate for P16)
    \item \textbf{Problem Prioritization}: Solve after securing easier early-mid band problems
    \item \textbf{Risk Assessment}: Medium risk - requires careful enumeration but concept is standard
\end{itemize}

\subsection*{B. Decision Points}
\begin{itemize}[leftmargin=*]
    \item \textbf{Soft Abandon}: If not seeing factorization connection after 2 minutes
    \item \textbf{Hard Abandon}: If arithmetic errors persist after 5 minutes
    \item \textbf{Strategic Skip}: If stronger in geometry, save this problem for later
\end{itemize}

\subsection*{C. Optimization Techniques}
\begin{itemize}[leftmargin=*]
    \item \textbf{Answer Choice Analysis}: Small answers (1-7) confirm exhaustive enumeration is expected
    \item \textbf{Estimation}: 690 has ~16 divisors, expect answer around 6-8
    \item \textbf{Partial Credit Note}: For COMC, showing the setup and a few valid sequences earns substantial credit
\end{itemize}
\end{competitionbox}

\begin{keyinsightbox}
\textbf{The Central Insight}: Transform the consecutive integer sum problem into a factorization problem. The equation $k(2n+k-1) = 690$ reveals that counting consecutive sum representations is equivalent to counting factor pairs of 690, subject to the constraints $k \geq 2$ and $n \geq 1$.

This insight converts a seemingly difficult counting problem into a straightforward divisor enumeration task.
\end{keyinsightbox}

\begin{warningbox}
\textbf{Most Common Pitfall}: Forgetting the constraint that $k \geq 2$ (at least two consecutive integers). The factor pair $(1, 690)$ satisfies the algebraic equation but violates the problem statement, leading to an incorrect answer of 8 instead of 7.

Always carefully check all constraints after finding algebraic solutions!
\end{warningbox}

\begin{tipbox}
\textbf{Pro Competition Tip}: When you see "consecutive integers" or "arithmetic sequence," immediately think of the formula $S = \frac{k(a_1 + a_k)}{2} = \frac{k(2n+k-1)}{2}$. This converts most such problems into algebraic manipulation or factorization tasks.

For "how many ways" counting problems with specific targets, prime factorization is often the key to systematic enumeration.
\end{tipbox}

\section*{Final Answer}
\boxed{\textbf{(E)}\ 7}

\section*{Confidence Assessment}
\begin{itemize}[leftmargin=*]
    \item \textbf{Confidence Level}: 95\% - Multiple verification methods confirm the answer
    \item \textbf{Verification Applied}: Direct calculation of all 7 sequences + factor pair counting
    \item \textbf{Flag for Review}: No - high confidence with thorough verification
    \item \textbf{Time Spent}: Approximately 3.5 minutes (within target range)
\end{itemize}

\end{document}